\chapter[Sermon 17. The Lord Commends the Virtues, Namely the Constancy of the Congregation of Philadelphia. Etc.]{The Lord Commends the Virtues, Namely the Constancy of the Congregation of Philadelphia. Etc.}
\label{sermon:017}
\markright{Sermon 17. The Lord Commends the Virtues, Namely the Constancy of ...}

And write unto the angel of the congregation of Philadelphia: this says he who is holy and true, who has the key of David, who opens and no one can shut, and who shuts and no one can open. I know your works. Behold, I have set before you an open door, and no one can shut it; for you have little strength. And you have kept my word, and you have not denied my name. Behold, I shall give some of the congregation of Satan, who call themselves Jews, and are not, but do lie. Behold, I will make them, that they shall come and worship before your feet.

\index{John, Saint}\index{Augustine, Saint}In all other congregations, the Lord at the least found some fault; in the only church of Philadelphia, he finds nothing to blame. Not that any person in this life is so perfect that he does not need God’s grace. For David cries out, “Enter not, Lord, into judgment with your servant, for no living person shall be justified in your sight.” But Saint John and Saint Paul also acknowledge all people to be subject to sin, a matter which Saint Augustine also discusses learnedly against the Pelagians. Therefore, when he finds nothing to blame in this congregation, it is not to be understood as though it were not tainted with daily faults, but rather that he imputes nothing because the sincerity and integrity of faith covers and hides whatever vice there may be. For there is no condemnation to those who are grafted in Christ Jesus. And although other churches also possess the true faith, this one excels especially, \&c. It might be referred chiefly to the Bishop of the same church.

In this sixth epistle, he commends the sincere faith and steadfastness of faith, and admonishes to persevere, proposing abundant rewards. And it contains much learning and diverse material, which will become apparent in the treatise itself.

And the Lord herein follows the same pattern, which we see he has followed in others. For it is one and the same doctrine with all churches and in all times. First, therefore, is shown to whom the epistle is written or dedicated: to the pastor and whole congregation of Philadelphia. Philadelphia was a city of Lydia, neither very famous nor yet obscure. We read how it has been often shaken by earthquakes and repaired again. Strabo mentions it in the 12th book of Geography, and so have other authors also. Yet it made itself famous by its virtues. After this, the Lord Christ is signified to be the Author of this epistle, who at other times also has told St. John what he should write.

\index{Antichrist}\index{Papacy}And to Christ are attributed three things, or rather Christ attributes three things to himself: that he is holy, true, and has the key of David. He has borrowed these from the image of the first chapter. Christ is holy because he is pure and clean from all defilement and all unrighteousness, very God, a consuming fire, doing no wrong to anyone, having nothing at all that may be blamed. For the Seraphim rightly say, “Holy, holy, holy, Lord God of hosts,” Isaiah 6:3. Christ is also the holiness of the saints, a sanctification that sanctifies all who are sanctified. He loves holiness in saints. Therefore, Christ is most truly called holy. Antichrist, the Pope, has taken upon himself this title, and sits on this beast so filthily, as if you were to call a privy or a jakes a rose garden. Spit upon that vile and filthy beast, which allows himself to be called the most holy father, and worship Christ, the holy one of all holy, unless you understand by that holiness not every holiness, but papal holiness—that is to say, stinking and swimming full of all abominations. Christ is likewise called true because he is eternal and faithful, evermore constant and incorruptible. He can neither deceive nor be deceived. He steadfastly keeps his promises. All his words are without doubt and true. Although flesh, which can endure no delay, often begins to doubt, yet not one point or jot of them falls away. The truth of the Lord endures forever. You stand upon a most sure foundation if you lean unto Christ, who in John 14 also calls himself the truth.

\index{Rome}Finally, he adds, “He has the key of David.” I spoke of the key in the first chapter. He alludes to chapter 22 of Isaiah, signifying the divine and almighty power of Christ, by which he brings us purified into the kingdom of heaven. This work, indeed, neither devils nor any power can hinder. It casts down the unholy into hell, and there is no one who can prevent it. Therefore, he says aptly and expressly, “He has,” not “He had” or “He will have,” but “He has now.” For he alone has this power, which he communicates to no other. The Pope of Rome lies, claiming to have this power. The only Son of God surpasses all others in this prerogative.

The apostles, as ministers and preachers, have received the keys of knowledge and of utterance, of learning, instruction, and introduction, by which they also, in threatening, exclude unbelievers from the kingdom of God and bind them in their sins. Almighty God, who has the ultimate authority, ratifies the judgment of the minister, which he pronounces not of himself, but according to Christ’s words. But these things agree well with what follows about the opened door, which no one can shut, and with the entire matter.

Now the Lord proceeds to tell what he wishes: and as he has said in all his letters, so he repeats in this also, that he knows all things concerning this and every other congregation.

And he praises the perseverance in faith within this congregation, indicating that it proceeds from the grace of Christ. “You have little power, and seemingly no strength, which the world considers power—riches, worldly wisdom, fortunate success, many friends, and things of that kind. Therefore, you can attribute nothing to yourself, nothing to your own strength, not even this: that you are a church, and that the truth of the gospel is freely preached among you.” For I have opened this door. And by my strength I keep it open, so that no one can shut this door—that is, to prohibit, hinder, or take away, by any means, the preaching and grace once granted. To open the door is a common phrase used by the Apostle in 1 Corinthians 16 and 2 Corinthians 2. He opens the door, which provides an opportunity and prepares the way to enter. Therefore, the word opened the door of life, so that the faithful might enter; the unbelievers could not stop this way. For the hand of Christ held the door open. And these things indeed declare, which explains why, in cities, towns, and villages not greatly furnished with power or strength, the course of the gospel proceeds with such fortunate success. And where many attempt to shut the door through schemes, trickery, threats, and persecutions, they cannot. These things are not accomplished through our skill and wisdom, but by the grace of God.

Nevertheless, if anyone desires to understand those things and what follows particularly concerning the pastor or bishop of the church, I will not object. For where he was humble and instructed with no worldly wisdom, yet furnished with God’s grace, he opened the way of salvation, which those who seek to abolish the preaching of the gospel now cannot shut. The power of Christ preserved him.

And now he more clearly urges or commends the faithful constancy in faith of the pastor and congregation: “You have kept my word, and have not denied my name.” When the Lord opened the door, lit the lamp, and gave spiritual gifts, the pastor and the congregation received them and kept them, and so did not deny them, nor trample them underfoot. This is excellent praise. Would God there were many such churches found today. Here also, O you church of Christ, you may learn all things—and each individual thing—about the duty of pastors, of the church, and of all godly men and women. You had no merit at all; God shone upon you through his grace. You had no worthiness, no deserving, no power nor authority: Christ has revealed himself to you through his mercy. Embrace him therefore who offers himself to you, hold fast, and never at any time let him go. \&c.

\index{Repentance}And note that the Lord says, my word, not every man’s word, but mine. What the word of Christ is, it is known to all men. For that which is written in the Gospel, and first indeed by the Prophets, and after by the Apostles was set forth in holy writ, is the word of Christ. It is not Christ's word that contends with this, although it be set forth by councils and holy fathers. Christ does not acknowledge that word; he acknowledges his own. And this must be observed and kept. The word of Christ is observed when it is not corrupted with additions, detractions, and wrappings, but is kept sincere in its natural sense. It is not kept when it is corrupted or depraved with men's inventions and perverse interpretations. The word of Christ is kept when it is commended not with the mouth alone, but is also expressed with godly works in the whole life, and beautified with holiness. It is not observed when, without repentance, men live most filthily. Finally, the word of Christ is observed and kept when it is not with any loathing of ours or impatience cast away, denied and forsaken. And therefore he affirms immediately, “You have not denied my name.” I have spoken elsewhere largely of confessing and denying of Christ's name. These things truly did the Philadelphians, and with these virtues through faith pleased the Lord. By these also may we commend ourselves to our Savior.

Furthermore, the Lord shows with how great a reward he would honor the constancy of the faithful in faith. “You now have many enemies,” he says, “because of your pure religion; but if you continue to hold fast, I will cause those same enemies to become your friends and ultimately companions in your faith. Those who have hitherto condemned you as evildoers and heretics will come to you with great humility, asking forgiveness, ready to receive your religion and to worship him whom they have blasphemed. And they shall come in the most humble manner, with the greatest humility possible. For so Isaiah foretold that it would come to pass in chapter 49, to which the Lord alludes at this present.

In the meantime, he addresses the Jews, singular enemies of the faith, whom he calls the Synagogue of Satan. For their teacher was none other than the Devil, as indeed they still have no better today. He calls them false Jews and liars. For neither did they confess the Lord, nor glorify God, nor believe in Christ their Messiah. But those who are Jews in deed are not such, as the Apostle Saint Paul said in the second chapter to the Romans. The power of God compelled many of them, forsaking their Jewishness, to embrace the Christian religion.

Therefore, if we desire or seek to retain also in our churches the pure word of God, and to receive our enemies humbly, we shall not achieve these things through wars or injustices, through railing and reproachful words, but by steadfast faith. But if either we do not profess our faith purely, or do not adorn it with virtues, what is it to be wondered at, though enemies remain enemies still, and continue to hate us every day more intensely than before, and at length oppress us, and extinguish the light of God’s word with many? Let us learn, dear brethren, by godliness, constancy, and holiness to win our brethren. The Lord Jesus grant us his grace to accomplish the same.
